\chapter{Surfaces elliptiques}

\section{Préliminaires sur les surfaces complexes}

\subsection{Définitions et premières propriétés}

Une \emph{surface complexe} est une variété analytique complexe de dimension 2, (donc une variété différentiable de dimension réelle 4). Plus précisément, il s'agit d'un espace topologique séparé $S$ muni d'un atlas de cartes à valeurs dans $\C^2$ dont les fonctions de transition sont holomorphes.

À toute surface complexe $S$ sont associés des invariants fondamentaux. Le \emph{genre géométrique} $p_g(S)$ est défini comme la dimension de l'espace des formes holomorphes de degré 2 :
\[
p_g(S) = \dim H^0(S, K_S),
\]
où $K_S = \Omega_S^2$ désigne le fibré canonique. Le \emph{genre irrégulier} $q(S)$ est défini par $q(S) = \dim H^1(S, \mathcal{O}_S)$. Ces invariants proviennent de la décomposition de Hodge de la cohomologie de $S$.

Un invariant fondamental en classification est la \emph{dimension de Kodaira} $\kappa(S)$, définie comme la dimension de l'image de l'application pluricanonique pour $m$ suffisamment grand, ou $-\infty$ si tous les espaces $H^0(S, mK_S)$ sont nuls. La dimension de Kodaira prend les valeurs $-\infty, 0, 1, 2$ et partitionne les surfaces en quatre grandes classes.

\subsection{Formule de Noether et caractéristiques d'Euler}

Pour une surface complexe projective, on dispose de la \emph{formule de Noether} :
\[
12 \chi(\mathcal{O}_S) = c_1(S)^2 + c_2(S),
\]
où $\chi(\mathcal{O}_S) = 1 - q(S) + p_g(S)$ est la caractéristique d'Euler-Poincaré du faisceau structural, $c_1(S)^2$ est le nombre d'auto-intersection du fibré canonique, et $c_2(S)$ est la caractéristique d'Euler topologique de $S$.

Cette formule, qui relie des invariants analytiques et topologiques, sera constamment utilisée dans l'étude des surfaces elliptiques.

\section{Définition des surfaces elliptiques}

\subsection{Définition et premières conséquences}

\begin{definition}
Une \emph{surface elliptique} est un couple $(S, \pi)$ formé d'une surface complexe $S$ et d'un morphisme surjectif $\pi: S \to C$ vers une courbe complexe $C$ (nécessairement lisse si $S$ est lisse) tel que :
\begin{enumerate}
\item La fibre générique $\pi^{-1}(t)$ est une courbe elliptique (surface de Riemann de genre 1) ;
\item La fibration est relativement minimale, c'est-à-dire qu'aucune fibre ne contient de courbe rationnelle lisse d'auto-intersection $-1$ (courbe exceptionnelle de la première espèce).
\end{enumerate}
\end{definition}

La condition de minimalité relative assure que l'on a éliminé toutes les contractions possibles de courbes exceptionnelles contenues dans les fibres. Elle est naturelle si l'on s'intéresse à la structure de la fibration elle-même.

\begin{remarque}
Une surface elliptique peut avoir des fibres singulières. L'ensemble des points $t \in C$ pour lesquels la fibre $\pi^{-1}(t)$ est singulière est fini.
\end{remarque}

\subsection{Invariants des surfaces elliptiques}

Soit $\pi: S \to C$ une surface elliptique. Notons $g$ le genre de la courbe $C$. On a les relations suivantes entre les invariants :

\begin{theoreme}
Pour une surface elliptique relativement minimale, on a :
\[
\chi(\mathcal{O}_S) = \sum_{t \in C} \delta_t,
\]
où $\delta_t$ est une contribution locale dépendant du type de la fibre singulière en $t$ (nulle pour une fibre lisse). De plus,
\[
c_2(S) = \sum_{t \in C} c_2(F_t),
\]
où $c_2(F_t)$ est la caractéristique d'Euler de la fibre $F_t$.
\end{theoreme}

Ces formules montrent comment les singularités des fibres contribuent aux invariants globaux de la surface.

\section{Classification de Kodaira des fibres singulières}

\subsection{Approche locale}

Soit $\pi: S \to C$ une surface elliptique et soit $t_0 \in C$ un point où la fibre $F_{t_0} = \pi^{-1}(t_0)$ est singulière. Au voisinage de $F_{t_0}$, on peut étudier la structure locale de la fibration. Kodaira a classifié toutes les configurations possibles de composantes irréductibles d'une fibre singulière, ainsi que leurs multiplicités et leurs intersections.

\begin{definition}
Soit $F = \sum_i m_i \Theta_i$ la décomposition d'une fibre singulière en composantes irréductibles. Le \emph{graphe dual} de $F$ est le graphe dont les sommets correspondent aux composantes $\Theta_i$ et dont les arêtes entre $\Theta_i$ et $\Theta_j$ comptent le nombre d'intersections $\Theta_i \cdot \Theta_j$.
\end{definition}

La classification de Kodaira repose sur l'étude de ce graphe dual et sur l'action de la monodromie locale autour du point singulier.

\subsection{Classification des types de fibres}

On distingue deux grandes familles de fibres singulières.

\subsubsection{Types $I_n$ ($n \ge 1$)}

Les fibres de type $I_n$ sont des cycles de $n$ courbes rationnelles lisses se coupant transversalement. Pour $n=1$, on obtient une courbe rationnelle avec un point double ordinaire (un nœud). Pour $n \ge 2$, on a un cycle de $n$ courbes $\Theta_1, \dots, \Theta_n$ avec $\Theta_i \cdot \Theta_{i+1} = 1$ (indices modulo $n$). La caractéristique d'Euler d'une fibre $I_n$ est $c_2 = n$.

\subsubsection{Types $I_n^*$ ($n \ge 0$)}

Les fibres de type $I_n^*$ sont des configurations plus complexes, obtenues à partir des types $I_n$ par adjonction d'une courbe supplémentaire rencontrant le cycle. Pour $n=0$, on obtient une configuration de 5 courbes rationnelles correspondant au diagramme de Dynkin $\tilde D_4$. Pour $n \ge 1$, on a des configurations de type $\tilde D_{4+n}$. La caractéristique d'Euler est $c_2 = n+4$.

\subsubsection{Types $II$, $III$, $IV$}

Ces types correspondent à des courbes rationnelles avec des singularités plus compliquées :
\begin{itemize}
\item Type $II$ : une courbe rationnelle avec un point de rebroussement (cusp). $c_2 = 2$.
\item Type $III$ : deux courbes rationnelles se touchant en un point (tangence). $c_2 = 3$.
\item Type $IV$ : trois courbes rationnelles concourantes. $c_2 = 4$.
\end{itemize}

\subsubsection{Types $II^*$, $III^*$, $IV^*$}

Ce sont les types duaux des précédents, obtenus par adjonction de courbes supplémentaires. Ils correspondent aux diagrammes de Dynkin $\tilde E_8$, $\tilde E_7$, $\tilde E_6$ respectivement. Leurs caractéristiques d'Euler sont respectivement $10$, $9$ et $8$.

\begin{table}[h]
\centering
\begin{tabular}{|c|c|c|}
\hline
\textbf{Type} & \textbf{Description} & $c_2$ \\
\hline
$I_0$ & Elliptique lisse & 0 \\
$I_n$ ($n \ge 1$) & Cycle de $n$ rationnelles & $n$ \\
$I_n^*$ ($n \ge 0$) & Configuration $\tilde D_{4+n}$ & $n+4$ \\
$II$ & Rationnelle avec cusp & 2 \\
$III$ & Deux rationnelles tangentes & 3 \\
$IV$ & Trois rationnelles concourantes & 4 \\
$II^*$ & $\tilde E_8$ & 10 \\
$III^*$ & $\tilde E_7$ & 9 \\
$IV^*$ & $\tilde E_6$ & 8 \\
\hline
\end{tabular}
\caption{Classification de Kodaira des fibres singulières}
\end{table}

\subsection{Propriétés des fibres singulières}

\begin{theoreme}[Kodaira \cite{Kodaira63}]
La liste précédente est exhaustive : toute fibre singulière d'une fibration elliptique relativement minimale est de l'un des types ci-dessus.
\end{theoreme}

Chaque type de fibre a une contribution spécifique à la caractéristique d'Euler de la surface :
\[
c_2(S) = \sum_{t \in C} c_2(F_t),
\]
et à la signature $\sigma(S)$ via la formule :
\[
\sigma(S) = -\frac{2}{3} \sum_t \delta_t,
\]
où $\delta_t$ est une contribution locale (par exemple, $\delta_t = 0$ pour $I_n$, $\delta_t = 1$ pour $I_n^*$, etc.).

\section{Exemples fondamentaux}

\subsection{La surface $E(1)$ (surface elliptique rationnelle)}

Considérons le plan projectif complexe $\CP^2$ et deux cubiques génériques $C_1$ et $C_2$. Par le théorème de Bézout, elles se coupent en 9 points distincts $p_1, \dots, p_9$. Soit $X = \Bl_{p_1, \dots, p_9}(\CP^2)$ l'éclatement de $\CP^2$ en ces 9 points. Le pinceau engendré par $C_1$ et $C_2$ définit une application
\[
\pi: X \longrightarrow \CP^1
\]
dont la fibre générique est une cubique lisse, donc une courbe elliptique.

Pour un choix générique des cubiques, la fibration présente exactement 12 fibres singulières, toutes de type $I_1$. La caractéristique d'Euler de $X$ est $c_2(X) = 12$, ce qui est cohérent avec la somme des $c_2$ des fibres. Cette surface, notée $E(1)$, est une surface rationnelle (birationnellement équivalente à $\CP^2$).

\subsection{Les surfaces K3 elliptiques}

Les surfaces K3 sont des surfaces simplement connexes avec $p_g = 1$ et $c_1 = 0$. De nombreuses surfaces K3 admettent des fibrations elliptiques. Un exemple classique est fourni par le complété lisse d'une quartique elliptique dans $\CP^3$, ou par le revêtement double du plan projectif ramifié le long d'une sextique.

Les surfaces K3 elliptiques ont généralement 24 fibres singulières (comptées avec multiplicités), et leurs types possibles sont contraints par la condition $c_2 = 24$.

\begin{exemple}[Surface K3 de Dolgachev \cite{Dolgačev1973}]
Dolgachev a étudié des familles de surfaces K3 algébriques contenant des courbes hyperelliptiques et a montré leur lien avec les surfaces elliptiques. Ces surfaces jouent un rôle important en symétrie miroir [citation:7].
\end{exemple}

\subsection{Les surfaces $E(n)$}

À partir de $E(1)$, on peut construire par fibre produit une suite de surfaces elliptiques $E(n)$ pour tout $n \ge 1$. La surface $E(n)$ est une fibration elliptique sur $\CP^1$ dont la monodromie satisfait $(UV)^{6n} = I$ dans $SL(2,\Z)$. Pour $n=2$, on obtient une surface K3 (car $c_2 = 24$). Pour $n \ge 3$, on obtient des surfaces de type général.

Les surfaces $E(n)$ sont essentielles en topologie différentielle : elles fournissent des exemples de 4-variétés simplement connexes avec des signatures arbitrairement grandes, et ont été utilisées par Gompf pour construire des variétés symplectiques exotiques.

\subsection{Les surfaces de Dolgachev}

Les surfaces de Dolgachev $E(1)_{p,q}$ sont des déformations de $E(1)$ obtenues par des \emph{twists logarithmiques}. Ce sont des surfaces elliptiques sur $\CP^1$ avec deux fibres multiples de types $I_0^*$ (ou plus généralement $I_n^*$) et de multiplicités $p$ et $q$. Historiquement, ces surfaces ont fourni les premiers exemples de surfaces homéomorphes (car simplement connexes avec mêmes invariants) mais non difféomorphes, jouant ainsi un rôle crucial dans le développement de la théorie de Donaldson.

\section{Formules invariantes}

Pour une surface elliptique $\pi: S \to C$, on a les relations suivantes reliant les invariants :

\begin{theoreme}[Formule de Noether pour surfaces elliptiques]
Soit $\pi: S \to C$ une surface elliptique relativement minimale. Alors :
\[
12 \chi(\mathcal{O}_S) = \sum_{t \in C} e_t,
\]
où $e_t$ est la contribution locale à la caractéristique d'Euler (égale à $c_2(F_t)$) et la somme porte sur toutes les fibres (y compris les fibres lisses, qui contribuent $0$).
\end{theoreme}

De plus, la signature $\sigma(S)$ est donnée par :
\[
\sigma(S) = -\frac{2}{3} \sum_t \delta_t,
\]
où $\delta_t$ vaut $0$ pour les types $I_n$, $II$, $III$, $IV$ ; $1$ pour $I_n^*$ ; $2$ pour $IV^*$ ; $3$ pour $III^*$ ; $4$ pour $II^*$ [citation:5].

Ces formules montrent comment les fibres singulières déterminent complètement les invariants topologiques fondamentaux de la surface.
